\section{Problem Statement}
\label{sec:problem}

\begin{figure*}[!t]
\centering
\includegraphics[width=\textwidth]{figures/research_object.pdf}
\caption{Hybrid AC/DC distribution system and transition-aware reliability
assessment. (a) Local VSC controls continuously determine converter currents,
whereas protection functions measure their assigned zones and trip only the
associated breakers or converters. The distribution-level controller receives
the isolated topology and available-resource status after the DAE simulation.
(b) For each joint $N$--$k$ contingency and operating state, the pre-restoration
DAE determines the VSCs that are unavailable after fault clearing. The
restoration MILP is then executed once using the initiating outages, the
DAE-derived VSC outages, and the original operating-state load and generation.
(c) Conventional assessment prescribes the restoration inputs. The proposed
assessment obtains them from protection--IBR interaction and reuses only the
restoration-entry state in subsequent reliability samples.}
\label{fig:research_object}
\end{figure*}

\subsection{Hybrid AC/DC Distribution Systems with Hierarchy Controllers}
\label{sec:research_object}
The assessed system is a hybrid AC/DC distribution network
$G=(\mathcal N^{\rm ac}\cup\mathcal N^{\rm dc},\mathcal E^{\rm ac}\cup
\mathcal E^{\rm dc}\cup\mathcal V)$, where $\mathcal N^{\rm ac}$ and
$\mathcal N^{\rm dc}$ are the AC and DC buses, $\mathcal E^{\rm ac}$ and
$\mathcal E^{\rm dc}$ are the corresponding branches, and $\mathcal V$ is the
set of interlinking VSCs. The system also contains loads $\mathcal D$,
converter-interfaced resources $\mathcal R$, controllable switches
$\mathcal E^{\rm sw}$, and protection devices $\mathcal Z$. The set of
physical components for which initiating outages are studied is denoted by
$\mathcal U\subseteq\mathcal E^{\rm ac}\cup\mathcal E^{\rm dc}\cup
\mathcal V$. As shown in Fig.~\ref{fig:research_object}(a), the physical
network is coordinated by local converter control, protection, and a
distribution-level restoration controller.

For each VSC or converter-interfaced resource $v$, the local controller uses
its terminal voltage, current, active power, and DC-link measurements to form
the current reference and operating mode. Current limitation, active/reactive
current priority, and fault ride-through (FRT) are included in this layer.
Therefore, local VSC control acts continuously during the fault and changes the
AC/DC voltage and current response seen by the protection devices. It does not
directly determine the post-fault switching plan.

Each protection function $z\in\mathcal Z$ is assigned a measurement set
$\mathcal M_z$, a protection zone $\mathcal P_z$, and a trip set
$\mathcal B_z$. Its pickup and definite-time states are driven only by the
available relay measurements. Once its operating condition is satisfied, the
corresponding breaker or converter in $\mathcal B_z$ is opened and the AC/DC
network equations are reinitialized. Hence, converter control and protection
are coupled through the same bus voltages and branch currents, although the two
controllers do not exchange their internal states. Complete observation of all
physical and controller states is not assumed.

The distribution-level controller is separated from the preceding transient
response. After fault isolation and completion of the DAE simulation, it
receives the isolated components, breaker positions, and reported VSC and DER
availability, and then solves the network reconfiguration and load-restoration
problem once. The resulting switch operations and power schedules are not fed
back into the DAE trajectory considered in this paper. Thus, local converter
and protection controls determine the actual restoration input, while the
system-level controller determines the load that can be supplied from this
input under the AC/DC network constraints.

\subsection{Reliability Assessment Scheme and Assumptions}
\label{sec:assessment_scheme}
Figure~\ref{fig:research_object}(b) gives the complete calculation procedure.
For an initiating contingency $k$ and operating state $\omega$, a consistent
pre-fault power-flow solution is first obtained. The initiating fault is then
applied to the coupled network, VSC-control, FRT, and protection DAE. When the
fault is cleared and all modeled protection actions are completed, the DAE
provides the set of unavailable VSCs at the restoration instant. This set,
together with the initiating outages, is supplied to the conventional
restoration MILP. The original load and generation of $\omega$ are retained,
and active and reactive load shedding are optimized for that operating state.
The annual reliability indices are finally obtained from the contingency
frequency, operating-state probability, restoration duration, and curtailed
load. The proposed method therefore changes the restoration-entry information,
but does not replace the conventional restoration model.

Four assumptions are made in this paper. First, $N$--1, $N$--2, and $N$--3
contingencies are represented as explicitly authored simultaneous joint events;
their frequencies are input data and are not calculated by multiplying
marginal component failure rates. Second, the pre-restoration model is a
balanced positive-sequence, phasor-domain mass-matrix DAE with average-value
VSCs, explicit current limits, FRT, and definite-time protection; traveling
waves, breaker arcs, and switching harmonics are outside the present scope.
Third, the DAE is terminated at a specified restoration instant after fault
isolation, and the system-level restoration is executed once without feedback
to the preceding trajectory. Fourth, a failed DAE initialization, event-time
network solution, or restoration optimization is retained as unresolved; no
normal restoration state or zero load shedding is assigned to an unresolved
condition.

\input{section_iii}
\input{section_iv}
